\section{Counting equivalence and homogeneity}
\label{sec:base-counting}

The central subgroup \(K=\ker\partial\) acts by adding a divergence-free
edge vector to the last coordinates of the atoms.  We now prove two facts
about this action.  First, a bounded number of variables cannot distinguish
the two charge assignments.  Second, within either structure, equality of
sufficiently wide counting types determines a \(K\)-orbit.  The second fact
will let us transport the supports of hereditarily finite sets.

\subsection{An expansion for the analysis}

Recall that an edge atom has coordinates \((e,u,z)\), whereas a configuration
atom at \(v\) specifies the states on all incident edges, subject to the face
compatibility equations and the divergence equation with right side \(b_v\).
Expand \(A_b\) to a structure \(\bar A_b\) as follows.  Partition its universe
into classes indexed by
\[
 (e,u)\quad\text{and}\quad(v,(u_f)_{f\in F(v)}).
\]
Here \(u_f\) is the common face coordinate on the two edges of \(f\) at
\(v\). Thus an edge class leaves only \(z\) unspecified, and a configuration
class leaves all incident \(z\)-coordinates unspecified.  Choose a linear order of
these class labels, the same for all charge assignments, and let
\(\preceq\) be the resulting total preorder on atoms.  Each edge class has
size \(3\), and a configuration class at \(v\) has size
\(3^{d(v)-1}\); in particular, corresponding classes are nonempty and have
the same size.  Finally add, for \(\delta\in\mathbb F_3\), the relation
\[
 \mathsf D_\delta((e,u,z),(e',u',z'))
 \quad\Longleftrightarrow\quad e=e'\ \text{and}\ z'-z=\delta,
\]
which is false on every other pair.  The vectors \(u,u'\) need not agree in
this definition. The preorder will let counting formulas identify a block
and its face coordinates. The relations \(\mathsf D_\delta\) will compare
scalar differences on the same edge, including between different face patterns.

Every element of \(K\) preserves the expansion.  The larger group \(H\)
need not preserve it.  Crucially, the preorder and the relations
\(\mathsf D_\delta\) are tools for this proof and are not part of the input
to a choiceless program.

Write \(C^h\) for first-order logic with equality, ordinary quantifiers and
exact counting quantifiers \(\exists^{=a}x\), using at most \(h\) variable
names.  Here \(a\) may be any nonnegative integer.  Formula length and these
integer constants need not be uniform in \(L\).  The \(C^h\)-type of an
\(r\)-tuple, for \(r\le h\), records all such formulas with free variables
among its \(r\) designated positions.  Coordinates of a tuple may repeat.

\subsection{A component that survives edge deletions}

Let \(G\) be the undirected graph underlying the oriented box grid on
\(V=\{0,\ldots,L-1\}^3\).  For \(R\subseteq V\), write
\(\delta_G R\) for its edge boundary; this notation distinguishes the
boundary of a vertex set from the incidence operator \(\partial\).
The projection estimate below is the three-dimensional discrete
Loomis--Whitney inequality \cite{loomisWhitney1949}.
We include its short proof and derive the needed finite-box boundary bound.

\begin{lemma}[Boundary of a box subset]
\label{lem:box-boundary}
Set \(c=1-2^{-1/3}>0\).  If \(\varnothing\ne R\subseteq V\) and
\(|R|\le L^3/2\), then
\[
 |\delta_G R|\ge c|R|^{2/3}.
\]
\end{lemma}
\begin{proof}
Let \(a_{ij},b_{ik},d_{jk}\) be the indicators of the three coordinate-pair
projections of \(R\).  If their sizes are \(A,B,D\), respectively, then
\[
 |R|\le\sum_{i,j,k}a_{ij}b_{ik}d_{jk}
 \le A^{1/2}\left(\sum_{i,j}
                  \left(\sum_k b_{ik}d_{jk}\right)^2\right)^{1/2}
 \le(ABD)^{1/2}.
\]
For the last inequality apply Cauchy--Schwarz to each inner sum and then
sum the product
\((\sum_k b_{ik}^2)(\sum_k d_{jk}^2)\) over \(i,j\).
Consequently at least one projection has size at least \(|R|^{2/3}\).
It indexes that many nonempty coordinate lines meeting \(R\).  At most
\(|R|/L\) of these lines lie wholly in \(R\).  Every remaining line is a
path meeting both \(R\) and its complement, so contains a boundary edge;
different lines give different edges.  Since
\[
 \frac{|R|}{L}\le 2^{-1/3}|R|^{2/3},
\]
the stated estimate follows.
\end{proof}

Set
\[
 M=\lfloor L^{7/4}\rfloor.
\]
For all sufficiently large \(L\), deleting any set \(T\) of at most
\(6M\) edges leaves a unique component containing more than \(L^3/2\)
vertices.  Call it the \emph{giant component}.  Indeed, summing the preceding
lemma over all components \(R\) of size at most \(L^3/2\) gives
\begin{equation}
\label{eq:small-components}
 \sum_R |R|
 \le c^{-3/2}\sum_R|\delta_G R|^{3/2}
 \le c^{-3/2}(2|T|)^{3/2}
 \le c^{-3/2}(12M)^{3/2}=o(L^3).
\end{equation}
The middle inequality uses \(\sum_R|\delta_G R|\le2|T|\) and
\(\sum_i a_i^{3/2}\le(\sum_i a_i)^{3/2}\) for nonnegative \(a_i\).
If every component had size at most half, the left side would equal
\(L^3\), a contradiction.  Two components cannot both have more than
half the vertices.  Moreover, whenever more edges are deleted within this
range, the new giant is contained in the old one: each new component is
contained in an old component, and its size identifies that old component
as the giant.

\subsection{Moving the unmatched charge}

The counting argument follows the finite-variable counting-game method
of Immerman and Lander \cite[Counting Quantifiers]{immermanLander1990}
and the charge-moving strategy of the CFI construction
\cite[Lemma~6.2 and proof of Theorem~6.4]{cfi1992}.
Here we construct the required bijections and prove formula preservation
directly; no game characterization is needed as a premise.

\begin{proposition}[Base counting equivalence]
\label{prop:base-equivalence}
For all sufficiently large \(L\), the expansions
\(\bar A_{b^0}\) and \(\bar A_{b^1}\), where
\(b^0=0\) and \(b^1=\mathbf1_{v_*}\), satisfy the same \(C^M\) sentences.
\end{proposition}
\begin{proof}
An assigned edge atom touches its edge, and an assigned configuration atom
at \(v\) touches all edges incident with \(v\).  At most \(M\) assigned
positions therefore touch a set \(T\) of at most \(6M\) edges.
Call a paired assignment \emph{compatible} if some vector
\(t\in\mathbb F_3^E\) and some vertex \(w\) in the giant of \(G\setminus T\)
satisfy
\begin{equation}
\label{eq:defect-flow}
 \partial t=b^1-b^0-\mathbf1_w,
\end{equation}
and each assigned atom is paired with the atom obtained by adding \(t_e\)
to its last coordinate on every relevant edge.  This operation leaves its
block and all face-vector coordinates unchanged.  The empty assignment is
compatible with \(t=0\) and \(w=v_*\).

A compatible assignment preserves every atomic formula.  On a fixed edge,
common addition of \(t_e\) preserves both raw differences and the corrected
differences defining \(\mathsf Z_\delta\).  It also preserves whether an
edge state is the specified neighbor of a configuration.  The block
relations, preorder, diagonal predicates and equality are preserved as
well.  This applies even when assigned coordinates repeat.

Deleting an assigned position preserves compatibility: the set of touched
edges shrinks, and the old giant is contained in the new one.  Suppose now
that at most \(M-1\) positions remain.  We construct one bijection of the
entire universes that extends them, regardless of which atom is next
assigned. Its definition may use a different shift on each block; every
such shift will agree with \(t\) on the retained touched edges.
Partition the universe into whole edge blocks and whole vertex
configuration blocks.  For such a block \(J\), let \(T_J\) be \(T\) together
with its edge, or with its vertex star, as appropriate.  There are at most
\(6M\) edges in \(T_J\).  Choose \(w_J\) in the giant of \(G\setminus T_J\).
Both \(w\) and \(w_J\) lie in the giant of \(G\setminus T\), so they are
joined by a path there.  A signed path vector \(p_J\), supported outside
\(T\), can be chosen with
\[
 \partial p_J=\mathbf1_w-\mathbf1_{w_J}.
\]
Put \(t_J=t+p_J\).  It agrees with \(t\) on every previously touched edge
and satisfies \eqref{eq:defect-flow} with \(w_J\) in place of \(w\).

Map every atom in \(J\) by its \(t_J\)-shift.  This gives a bijection on an
edge block.  For a configuration block at \(v\), deleting the star isolates
\(v\), so \(v\ne w_J\).  Hence
\((\partial t_J)_v=b^1_v-b^0_v\), exactly the condition for translation to
map its configurations bijectively to those in the other structure; the
face-vector constraints are unchanged.  The union of these block
bijections is a bijection of universes. For a new atom in block \(J\),
both it and all retained positions are paired by the single shift \(t_J\).
Thus adjoining that atom extends the compatible assignment, with witness
\(t_J,w_J\).

Induct on formulas.  Atomic formulas were handled above, and Boolean
operations preserve agreement.  At a quantifier, first discard any old
assignment of its variable, leaving at most \(M-1\) positions.  The
bijection just constructed pairs all extensions on which the body has the
same truth value by induction.  It therefore preserves ordinary
quantification and every exact counting quantifier, including count zero.
Starting from the empty assignment proves the proposition.
\end{proof}

\subsection{Recognizing the central-group orbits}

For the next step, the two tuples lie in the \emph{same} expansion.  A shift
prescribed on their touched edges must extend to a divergence-free vector.
The only obstruction is a nonzero net prescribed flow across one of the
remaining components.

\begin{lemma}[Extending a partial flow]
\label{lem:partial-flow}
Let \(T\subseteq E\) and \(\sigma\in\mathbb F_3^T\).  There is
\(k\in\ker\partial\) with \(k|_T=\sigma\) if and only if, for every component
\(U\) of \(G\setminus T\),
\begin{equation}
\label{eq:component-balance}
 R_U:=\sum_{v\in U}\sum_{\substack{e\in T\\e\ni v}}
                      \epsilon_{ve}\sigma_e=0.
\end{equation}
If an extension does not exist, there is a nonempty failing component
with
\[
 |U|\le L^3/2,\qquad |U|\le(|T|/c)^{3/2}.
\]
\end{lemma}
\begin{proof}
Extend \(\sigma\) by zero outside \(T\).  The required values on \(E\setminus
T\) must have divergence \(-\partial\sigma\).  On any connected graph, the
image of its incidence operator is exactly the vertex vectors with zero
sum.  To see sufficiency without a rank calculation, take a spanning tree,
set nontree edge values to zero, and eliminate leaves: the divergence at a
leaf uniquely determines the value on its tree edge.  The prescribed total
sum makes the final root equation hold.  Apply this separately to each
component, allowing an isolated vertex as a one-vertex tree.  The necessary
zero-sum conditions are precisely \eqref{eq:component-balance}.

The sum of all \(R_U\) is zero, because the two incidence signs of each edge
cancel.  Thus failure at one component implies failure at at least two, so
one failing component has size at most \(L^3/2\).  Its boundary is contained
in \(T\), and Lemma~\ref{lem:box-boundary} gives the second estimate.
\end{proof}

\begin{proposition}[Quantitative homogeneity]
\label{prop:base-homogeneity}
Let \(r\ge1\), and suppose
\begin{equation}
\label{eq:homogeneity-width}
 7r+7(6r/c)^{3/2}+1\le M.
\end{equation}
If two \(r\)-tuples in the same \(\bar A_b\) have equal \(C^M\)-types, then
some \(k\in K\) maps the first tuple to the second.  The corresponding
statement for empty tuples holds without any width hypothesis.
\end{proposition}
\begin{proof}
Write the tuples as \(\bar a,\bar a'\).  Each preorder class is definable
using two variable names.  If \(h\) is the number of atoms in strictly
earlier classes, its predicate is
\[
 \theta_h(x):=\exists^{=h}y\,
       (y\preceq x\land\neg(x\preceq y)).
\]
Different classes give different \(h\), since every class is nonempty.
Type equality therefore puts corresponding tuple coordinates in the same
class.  They have the same blocks and face-vector patterns.

Let \(T\) be the touched edges, so \(|T|\le6r\).  An \emph{occurrence} of a
state on \(e\in T\) is either an edge atom appearing in a tuple coordinate
or the state on \(e\) in a configuration appearing there.  In the latter
case that state is the unique \(y\) satisfying \(\mathsf I(a_i,y)\) and its
specified edge-class predicate.  The corresponding occurrence for
\(\bar a'\) has the same edge and face-vector coordinates.  For any two
occurrences on one edge, their raw difference can be tested by
\(\mathsf D_\delta\), quantifying their uniquely specified neighbors if
necessary.  These tests use at most the \(r\) tuple names, two neighbor
names, and one name reused in class-counting predicates.  Condition
\eqref{eq:homogeneity-width} in particular permits these \(r+3\) names.
Type equality shows that every pair of corresponding occurrences has the
same raw difference.  Hence there is a single well-defined
\(\sigma_e\) equal to the change in \(z\) from \(\bar a\) to \(\bar a'\)
for all occurrences on \(e\).

Suppose \(\sigma\) cannot be extended to \(K\).  Choose a failing component
\(U\) from Lemma~\ref{lem:partial-flow}; then
\[
 |U|\le(6r/c)^{3/2}.
\]
We construct a formula true at \(\bar a\) and false at \(\bar a'\).
For every \(v\in U\), choose any configuration \(g_v^0\) in the first
assignment's structure and let \(h_{v,e}^0\) be its neighbor on every incident
edge. The chosen configurations need not agree across an internal edge.
The formula will preserve their original raw difference, so any new
witnesses will have equal scalar changes at its two ends. These changes
will cancel when the divergence equations are summed over \(U\).
On each boundary edge \(e\in\delta_GU\subseteq T\), choose one
occurrence from \(\bar a\), and denote its state by \(o_e^0\).  Write
\(o_e\) for the corresponding tuple variable, or the auxiliary neighbor
variable when one is needed.  The formula existentially quantifies
variables \(g_v,h_{v,e}\), and, if the chosen occurrence comes from a tuple configuration,
one neighbor variable \(y_e\).  Its conjunction asserts:
\begin{enumerate}
\item each \(g_v,h_{v,e}\) is in the class of the chosen atom, and
      \(\mathsf I(g_v,h_{v,e})\) holds;
\item for every edge internal to \(U\), the raw difference between its two
      variables \(h_{v,e}\) equals the difference between the two chosen
      atoms \(h_{v,e}^0\);
\item for every boundary edge, the raw difference between \(h_{v,e}\) and
      its chosen tuple occurrence equals their original raw difference;
      if that occurrence requires \(y_e\), also impose its edge-class
      predicate and its incidence with the relevant tuple configuration.
\end{enumerate}
For definiteness, on an internal edge joining \(v,w\), choose an ordering
of its endpoints and use
\[
 \mathsf D_{z(h_{w,e}^0)-z(h_{v,e}^0)}(h_{v,e},h_{w,e});
\]
on a boundary edge incident with \(v\in U\), use
\[
 \mathsf D_{z(h_{v,e}^0)-z(o_e^0)}(o_e,h_{v,e}),
\]
where both atoms in the subscript are the fixed original choices.
Thus \(o_e=y_e\) in the arguments when the occurrence is supplied by a
tuple configuration.
An edge in \(T\) with both endpoints in \(U\) is treated as internal; its
contribution to \(R_U\) cancels and requires no boundary link.
All class predicates are of the form \(\theta_h\) above, and all difference
conditions use one of the three \(\mathsf D_\delta\).  This describes a
finite formula in the expansion's vocabulary; the numerical constants and
its length may depend on the chosen assignments.  The indicated choices
witness its truth at \(\bar a\).

If it were true at \(\bar a'\), compare the realized \(z\)-coordinates of
its witnesses with those of \(h_{v,e}^0\), writing the changes as
\(d_{v,e}\).  Both configurations at \(v\) have divergence \(b_v\), so
\(\sum_{e\ni v}\epsilon_{ve}d_{v,e}=0\).  On an internal edge the difference
constraint gives equal changes at its two ends.  On a boundary edge it
forces \(d_{v,e}=\sigma_e\).  Summing the vertex equations over \(U\)
therefore gives \(R_U=0\), contrary to the choice of \(U\).

There are \(r\) tuple variables, \(|U|\) configuration variables, at most
\(6|U|\) incident-state variables, and at most \(|\delta_GU|\le6r\) extra
occurrence variables.  One further name can be reused in every class
predicate.  Thus the formula uses at most
\[
 r+7|U|+6r+1\le7r+7(6r/c)^{3/2}+1\le M
\]
names, contradicting type equality.  The partial shift therefore extends
to some \(k\in K\), whose action maps every coordinate of \(\bar a\) to the
corresponding coordinate of \(\bar a'\).  For empty tuples take \(k=0\).
\end{proof}

For the support scale used below, fix \(q,m\) first and put
\[
 s=\left\lceil2L(q\log_3 N+1)^2\right\rceil+1,
 \qquad N=\Theta(L^3).
\]
Uniformly for \(1\le r\le\max\{2,m\}s\), the left side of
\eqref{eq:homogeneity-width} is
\[
 O_{q,m}\bigl(L^{3/2}(1+\log L)^3\bigr)=o(L^{7/4}).
\]
Consequently Proposition~\ref{prop:base-homogeneity} applies throughout
that range when \(L\) is sufficiently large. The same choice of \(L\)
ensures \(\max\{4,m+1\}s\le M\), the separate width requirement for
extending atom types in the hereditary-set comparison below.
The exponent \(7/4\) simultaneously ensures \(M^{3/2}=o(L^3)\) for the
giant-component estimate and \(L^{3/2}(1+\log L)^3=o(M)\) for homogeneity
at the support scale. Any fixed exponent strictly between
\(3/2\) and \(2\) would meet these asymptotic requirements.
All these conclusions hold for fixed \(q,m\), chosen before \(L\).
